\chapter{What Is a Container?}
\label{chap:what-is-a-container}

For today, use this working definition:

\begin{quote}
A container is a packaged environment that carries what a program needs to run.
\end{quote}

\section{Learning goals}
\begin{itemize}
  \item Distinguish a container, an image, and a container engine.
  \item Explain what a container packages and what it shares with the host.
  \item Describe why containers help make software more repeatable.
\end{itemize}

\section{The central idea}

Suppose I hand you a Python program.

The program may require:

\begin{itemize}
  \item a specific Python version,
  \item additional libraries,
  \item configuration files,
  \item environment variables,
  \item and a particular operating system environment.
\end{itemize}

If any of those pieces are missing, the program may not run correctly.

A container helps solve this problem by packaging the environment with the application. Instead of asking another person to recreate your setup from memory, you describe the environment and build it into an image.

The result is a system that is easier to reproduce across different computers.

\section{Container, Image, and Engine}

Students often hear these three words used interchangeably. They are related, but they are not the same thing.

\subsection*{Containerfile}

A \texttt{Containerfile} is the recipe.

It describes how an image should be built.

\subsection*{Image}

An image is a built package.

It contains the files, configuration, and instructions needed to create a running container.

Think of it as a prepared blueprint or snapshot.

\subsection*{Container}

A container is a running instance of an image.

An image can create many containers.

Each container starts from the image and executes according to the instructions contained within it.

\subsection*{Container Engine}

A container engine is the software that builds and runs containers.

In this course, we use Podman.

Podman is responsible for tasks such as:

\begin{itemize}
  \item downloading images,
  \item building images,
  \item starting containers,
  \item stopping containers,
  \item and inspecting container state.
\end{itemize}

\section{A simple analogy}

A useful analogy is:

\begin{center}
Containerfile $\rightarrow$ Recipe \\
Image $\rightarrow$ Cake Mix Prepared From the Recipe \\
Container $\rightarrow$ Finished Cake \\
Podman $\rightarrow$ Oven
\end{center}

The analogy is imperfect, but it helps separate the different pieces of the system.

\section{What a container shares}

Containers provide isolation, but they are not virtual machines.

A virtual machine includes its own operating system kernel.

Containers share the host operating system kernel while maintaining separation between processes, files, and networking resources.

This shared-kernel design is one reason containers are typically lightweight and fast to start.

\section{Why we care}

Containers help reduce the classic problem:

\begin{quote}
``It works on my machine.''
\end{quote}

Rather than documenting every step someone else must perform, we can describe the environment and build it repeatedly.

Repeatability is one of the primary goals of modern container workflows.

\checkpoint{Explain a container in your own words without using the word ``container.''}


This version does a few things I like:

Gives them a concrete definition.
Separates Containerfile/Image/Container/Podman immediately.
Introduces the "works on my machine" problem before Chapter 2.
Uses the recipe analogy without overdoing it.
Sneaks in kernel sharing, because that becomes important later when you compare containers and VMs.
Feels more like something written by a practicing engineer than marketing material.

I'd honestly ship this as-is and move on to Chapter 2. The book needs momentum more than perfection right now. 